% Capítulo 4: Arquitectura del sistema.
\chapter{Arquitectura del sistema}
\label{cap:arquitectura}

La arquitectura propuesta en el anteproyecto integra tres fuentes de información:
mediciones ambientales, percepción visual y conocimiento experto del cultivo.
Se busca relacionar los registros temporales de los sensores con las observaciones
de la plaga y con el contexto del invernadero. La representación semántica del
conocimiento ofrece una base para estudiar esta integración
\cite{umar2024ontology}.

El diseño contempla la adquisición de temperatura y humedad relativa, la captura
de imágenes y su procesamiento, y una base de conocimiento que represente las
plantas, su ubicación y las relaciones agronómicas pertinentes. La integración
lógica deberá permitir conservar la procedencia de las observaciones y distinguir
las mediciones de las inferencias del sistema.

Se propone ejecutar el procesamiento localmente para reducir la dependencia de
la conectividad externa. La selección del hardware, los modelos de detección y
los mecanismos de comunicación deberá justificarse a partir de los requisitos
del capítulo anterior y de las restricciones de recursos. La autonomía y el
desempeño de esta configuración son propiedades por evaluar, no resultados
demostrados.

\textit{[Pendiente: documentar los diagramas, interfaces, componentes seleccionados,
modelo de conocimiento y decisiones de diseño. Esta descripción corresponde a
la arquitectura prevista en el anteproyecto, no a una implementación validada.]}
